\scanpage{13}
Sur \(]0,+\infty[\), \unsure{si|et} on dit qu'une fonction \(P\) sur \(L^\infty(]0,+\infty[)\) est une fonction de Polya si et seulement si pour \deleted{tout} \(0<\omega<\infty\), sa restriction à \(L^\infty(]0,\omega[)\) est une fonction de Polya.

\margin{Exemple de fonction de Polya : \(P(f)=\int_0^{\unsure{\omega|\infty}}\alpha(f)\,\unsure{dt|dm}\)\par où \(\alpha\) est une fonction convexe sur \(R\)\par (\unsure{crois. en \(f\) si} \(\alpha\) est croissante (on suppose \(\alpha(0)=0\) si \(\omega=+\infty\))}

\deleted{\unsure{Notons} \(\unsure{\|f\|=0}\)\quad \(\unsure{f\longmapsto f+\|f\|}\)}

\deleted{Soit \unsure{maintenant} \(\unsure{M|\omega}\) un \illegible{2} \unsure{réel}. \deleted{\unsure{toute}} \deleted{\unsure{par|pour}} fonction de Polya sur \(L^\infty(]0,\omega[)\) \unsure{(où} \(0<\omega<+\infty)\) est la restriction d'une fonction de Polya sur \(K^\infty(]0,+\infty[)\), à savoir celle}

\deleted{Soit \(M\) un espace mesuré, \(P\) une fonction de Polya}

\uline{Remarque} Fonction de Polya \unsure{non} \unsure{croissante}.
\inserted{\unsure{condition} \unsure{définie}.}

\uline{Définition 2} \uline{Fonction de Weyl} sur \(L^{\infty+}(]0,\omega[)\) est une fonction satisfaisant à : \(P_1,P_3,\) \deleted{\illegible{1}} \inserted{continue} pour la norme, et à
\deleted{\(\unsure{P'_3|P_3}\) \unsure{sur tout compact}}
\[
(W_2)\qquad P(\exp f)\quad\text{est une fonction convexe sur}
\]
\[
L^\infty(]0,\omega[).
\]

\figure{\begin{tikzpicture}\draw (0,0)--(0,4);\end{tikzpicture}}

Une fonction sur \((K^\infty)^+(]0,+\infty[)\) est une fonction de Weyl, si sa restriction à tout \((K^\infty)^+(]0,\omega[)\) \((0<\omega<\infty)\) l'est.

\unsure{Supposons|Donnons} \unsure{en fait|enfin}, \(W\) fonction de Weyl sur \(L^\infty(]0,\omega[)\). Alors \deleted{\(\unsure{P\circ}\)} \(W(\exp f)\) est une fonction de Polya sur \(L^\infty(]0,\omega[)\). Par suite, on a
\[
P(\exp f)=\sup_i\left(\unsure{c_i|a_i}+\int_0^\omega\varphi_i\varphi_{\deleted{\unsure{e^f|\exp}}f}\right)
\]
\unsure{soit}
\[
P(f)=\sup_i\left(\unsure{a_i|c_i}+\int_0^\infty\varphi_i\log\varphi_f\right)
\]
si \(f\) est une fonction \deleted{positive} telle que \deleted{\(\unsure{m(f)}\)} \inserted{\(\unsure{p(f):|P(f):}\)}
\[
\inf_{0\le s\le\omega}f(s)>0.
\]
La connaissance \inserted{de \(W(\exp f)\), i.e.} des \(\varphi_i,\unsure{c_i|a_i}\), détermine alors \(W\) entièrement, \deleted{\unsure{car}} grâce à la continuité pour la norme. Le cas le plus important est \unsure{celui où} \(\unsure{W}\) est croissante. Une fonction de Polya sur le

\mspage{14}
